\section{百篇精读后的证据链：Hall效应改变了什么}
本轮39篇按模型争议、可复现细节和独立实验价值选择，而非按篇幅或容易下载来凑数。
以下侧重把新论文放回已有证据链；逐篇页码、参数和限制见第七章。

\subsection{早期闭合与移除Hall对照：不能只比较开关}
\citet{L029}的Hall/混合比较采用特定经验电子压力张量闭合，且其方程省略磁旋转项；这与今天完整十矩不是同一模型。
它提示离子非对角压力输运能改变层结构，不能从全局磁通相近推断离子应力无关。
\citet{L154}对无Hall快速重联的检查更具体：粗格Sweet--Parker层约2.5个网格时出现岛，约5格时可恢复较稳定层；去掉Hall后的率仍约0.015，含Hall约0.03，但仍有超电阻等耗散。
因此该对照揭示分辨率与耗散的作用，不能单独证明所有系统中Hall既必要又充分。

\citet{L035}的边界压缩算例中，混合模型因离子回旋效应出现更宽的响应，进入快速阶段的时间与Hall模型不同；局部嵌入Hall区可复现其特定总通量演化。
这解决的是\emph{该驱动条件下局地两流体区域是否足以改变整体响应}，未解决任意装置边界下的模型耦合。

\subsection{电阻从公式中消去，仍有未选定的长度}
\citet{L145}用电阻拖曳同时决定电子排出和耗散层厚，得到$E\sim B_mV_{\!A,m}d_i/L_e$；电阻消去后$L_e$和$B_m$仍是外部输入。
\citet{L137}从给定Harris外场及给定电场构造弱Hall局地解，也未选择全局率，且局部展开出现中心O点等几何问题。
\citet{L160}的有限质量比稳态分支又显示：当电子层厚受到$d_e$限制，$E\sim\eta B_m/(\mu_0d_e)$可显含电阻。
这些研究并不互相取消，而是采用不同闭合、长度约束和稳态假设；“快率与微观项无关”必须附带具体分支条件。

\citet{L089}从稳态不可压EMHD出发，用单位磁通的体积积分解释四极Hall场与电子流函数的关系。
该推导最有价值的细节是：理想EMHD内区的面外场主要由电子输运，生成需联系离子过渡区；预设X形面内场还在分离线产生对数奇异，必须靠惯性等修正。
它解释四极结构，未独立确定重联率和电子层长。\citet{L119}的实验只在远离X点和分离线的部分区域支持其沿射线近似恒定的场幅，近区明显偏离。

\subsection{线性Hall加速有条件，不能当成理想Hall直接断线}
\citet{L123}研究线绑边界的线性X点衰减；弱Hall拟合同时依赖$d_i$和电阻，因此将尺寸变化只归到$d_i^2$会漏掉归一化电阻的变化。
\citet{L122}提出三层变分近似，却没有完整$\Delta'$外匹配，且高斯试探电流导致某个$1/k^2$积分需额外正则化说明。其弱Hall修正还存在可直接代数复核的系数问题。
故本报告保留其启发，不将其形式上有限的无电阻增长解释成平滑理想Hall中磁通已被打破。

更系统的Hall撕裂分析\citep{L250}区分可压缩、常磁通与非常磁通分支；某一分支增长不显含$S$，仍需有限电阻完成内层匹配。
\citet{L238}在受迫稳定Taylor平衡中找到尤其值得注意的中间区：早期已出现Hall动力学，后期却能回到MHD式受限演化，尽管片厚小于$d_i$。
Hall还可能提高单模增长却增大最快波长，使可容纳的磁岛更少；“更快增长”与“更多磁岛”不是同一句话。

\section{三维并不只有“更快”：减速、倾角与传播记忆}
\subsection{电子片展宽与Hall场减弱也属于新现象}
\citet{L020}的三维惯性双流体计算中，重联建立的梯度激发静电涨落/低混合漂移类活动，电子层从约$d_e$展宽到$2.5d_e$，率比对应二维低约25\%，Hall场幅降低约30\%--40\%。
其平均惯性关联更像动量输运或异常黏性，不能统一称作异常电阻；密度梯度的放松还与磁场输运有关。
该算例质量比约278，$d_e$约3格，盒第三方向仅$1.25d_i$，因此是有条件的减速证据，不是三维必然减速定律。

\citet{L204}则发现Hall驱动的场向电流空间图样可不同于MHD的射流偏转机制；即使电流片较厚，系统沿电流方向的尺度与片厚之比仍会控制哪类图样主导。
模型没有地球偶极场、电离层反馈或电子加速闭合，不能直接声称重现完整亚暴电流楔。

\subsection{斜模更快与二维全局最快可以同时成立}
\citet{L224}在强导场force-free平衡、电子惯性和较弱超电阻下，固定波数比较发现小倾角撕裂比平行模快；导场比约10时部分倾角可快约两倍。提高超电阻又使平行模重新占优。
\citet{L268}使用电阻、不可压Harris平衡，完整扫描$(k_x,k_z)$后发现：某些固定较大$k_x$切片上斜模略快，但若盒能容纳较长波，全局最大增长仍在$k_z=0$。
两者的平衡、电子非理想项、扫描维度都不同，不能抽掉这些条件后写成互相证伪。
后者的复频率还表示扰动沿片传播，线性相速度不等于重联前沿扩展速度。

\citet{L166}三维不对称算例测得X线倾角约$35.9^\circ\pm8.5^\circ$，同时涵盖最大出流准则约$34.3^\circ$和最大率准则约$31.7^\circ$。
因此这是对角度量级的支持，精度不足以在两理论中作唯一选择；局地电流判据、有限盒和电子尺度分辨率也限制精确判别。

\subsection{重联沿非均匀电流片扩展会保留上游历史}
\citet{L294}得到薄片向厚片扩展时
$v_s\sim\VA d_iw_1/w^2$，厚片向薄片时约为$\VA d_i/w$。
前一种含起始薄片厚度$w_1$，因此前沿速度并非只由当前局地片厚决定；薄处生成的重联磁通进入厚处后被稀释，是这一“记忆”的来源。
这是比均匀片双向/单向传播更具体的新预测，适合与局域启动实验及非均匀磁尾比较。

论文表中的个别差值符号和磁尾应用时间需要独立检查。
按文内$w_1=0.89$或$3.5d_i$、$w_2=8.9d_i$、过渡指数2、距离$133d_i$、$d_i=720$ km、$\VA=1400$ km/s积分，相应传播时间约10.0--23.5分钟；不直接沿用正文2.5分钟的下限。
这是对该组输入的复算，不是对真实亚暴时间作新拟合；文后正式勘误涉及的参考文献/图例问题也不能自动视为已修正这些数值。

\section{增长加速、磁岛、涡旋和压强抑制的竞争}
\subsection{爆发律的更新伴随计算分辨率的更新}
\citet{L191}以能量变分纠正冷惯性模型的非线性爆发指数，\citet{L214}在暖双场模型中强调局域X几何及$\dot\epsilon\propto\epsilon^{7/4}$。
\citet{F697}进一步比较离子/电子温度与弱电阻：固定岛幅时某些速率和能量释放曲线对微观尺度较弱敏感，但从种子开始的总时间仍依赖线性内层尺度。
更关键的是，其冷模型会出次级岛并间歇减速；作者明确指出早期横向模数少和数值黏性漏掉了磁岛。
因此新结果是“局域爆发与磁岛间歇相竞争”，而不是一条适用于所有无碰撞重联的幂律。

\subsection{EMHD岛内部还会发生另一种不稳定性}
\citet{E459}在冷惯性EMHD中发现，早期有序的电流片级联被岛内射流的Kelvin--Helmholtz类涡旋破坏；电子压缩性可延迟或抑制这一过程。
它考察的主要是磁岛内部电流/涡量重新组织，不应只按“撕裂产生更多磁岛”概括。
模拟最高另做到$2048\times512$；晚期能量/Casimir误差可达5\%--10\%，不能由衰减曲线直接推断真实电子加热。

\citet{L202}则在惯性/电阻EMHD线性问题中强调纯增长本征模与驻波whistler不同，并构造有限中心电流的解。
其与MRX的增长时间和宽度只是尺度估计：计算约1.2 cm宽、实验比较约3.3 cm，且从假定10\%初扰动估总时间；未直接模拟稳态MRX排出和全欧姆预算。

\citet{L230}的强导场X点塌缩结果来自PIC，应放在动理学对照类。
其中电子涡旋可伴随次级X点，部分区域由对流惯性支撑电场；三维出现倾斜涡管。
但其三维图没有相应完整欧姆预算或全局率诊断，$Q$涡旋指标也不是重联判据，不能据图证明“电子压力项在三维全处消失”。

\subsection{磁岛触发Hall之后，重联还可能因压强结构被关停}
\citet{L159}在一组大$S$计算中看到Sweet--Parker、磁岛增强、Hall快速阶段先后出现，Hall阶段排出竞争磁岛。
后续文献已有开放X点持续排岛的情况，因此“Hall与磁岛互斥”只能限于该组边界、尺度与时段。
\citet{L174}则把压强梯度带入重联层，发现即使还有可释放磁能，反向抗磁漂移相对出流足够大时也可抑制重联。
该模型强导场、冷离子、绝热电子压力，MAST比较又依赖尺度估计；它提出了饱和/关停候选机制，尚非完整托卡马克锯齿预测。

\citet{E472}研究的是另一个问题：初始密度梯度使线性岛传播并降低增长；热离子时传播可朝离子抗磁方向。
其解析频率/增长率与数值有可见偏差，并未等价地验证\citet{L174}的非线性出流竞争条件。

\section{回旋平均、能谱与加热：全局相似不意味着模型等价}
\citet{E471}的四场陀螺流体显示，热离子改变导心流、X点尖角、电子密度分离线结构和高谐波，而部分全局磁通仍相近。
其$N_i$是导心密度，不是仪器测到的粒子密度；把两者直接比较会把极化响应误判成电子--离子非准中性。
\citet{E474}加入有限$\beta_e$和电子FLR后出现非线性减速间隔与第二轮超指数阶段，电子密度出现嵌套四极及更平缓梯度，而广义不变量仍继续细丝化。
此结果的方向随固定参数组改变，见式\eqref{eq:selected-beta}；部分增长率振荡还来自X点移动。

\citet{L264}比较Hall与混合动理学湍流时，磁场谱相似而速度及能量分配可不同；对比时段对应不同成熟阶段，不能仅由频率相近建立完整因果能量通路。
\citet{F644}的KH驱动Hall谱在离子尺度附近变陡，并在小电流片识别重联候选；但模型没有显式电子破冻结项，数值滤波就是必须审计的部分。
这些研究支持Hall模型捕获部分多尺度磁结构，尚未证明它精确预测粒子热化。

\citet{E512}提供了一个清楚的边界：电子平行相混合可把流体自由能送到高Hermite矩，再由弱碰撞热化。
这需要流体--动理学扩展，不能用纯Hamiltonian双流体的磁能下降代替。
\citet{E553}又从分布重建说明有限矩的内容边界：即使21矩恢复三角形主体，仍缺双峰与精细多轨道群。
\citet{L313}所讨论的电子正则磁场冻结则强调惯性拓扑变化不必耗散能量；它是理论框架讨论，没有本篇新的大规模重联模拟。
三维指数分离可缩短特定流程中的尺度形成时间，但三维并不自动保证混沌，常密度固定离子前提也不可随意推广。
